\begin{recipe}
[% 
    preparationtime = {\unit[25]{min}},
    bakingtime = {\unit[5--8]{min}},
    bakingtemperature = {\protect\bakingtemperature{
        topbottomheat=\unit[200]{\textcelcius}
    }},
    portion = {\portion{2}},
    %source = {},
    calory = {240kcal | 3g Protein | 33g KH | 12g Fett},
    price = {1,60€},
    created = {04.06.2026},
    rating = {2},
]
{Reispapier-Pillows mit Guacamole}

\graph{% pictures
    big=pic/vorspeise/reispapier-pillows-guacamole
}

\introduction{%
    Mal etwas, das man richtig fancy anrichten kann, Geschmack so semi, aber sieht cool aus. Definitiv schnell essen, sonst wird das Reispapier wieder weich.
}

\ingredients{
    \textbf{Pillows} & \\
    2 Blätter & Reispapier\index{Getreide!Reispapier}\\
    & Wasser zum Anfeuchten\\
    \textbf{Guacamole} & \\
    1 & Reife Avocado\index{Obst!Avocado}\\
    \nicefrac{1}{2} & Limette\\
    Optional & Koriander oder Petersilie\\
    & Salz, Pfeffer\\
    & Chiliflocken
}

\preparation{%
    \step%
    Backofen auf 200 °C Ober-/Unterhitze vorheizen. Ein Backblech mit Backpapier belegen oder ein Gitter auf ein Backblech setzen.
    
    \step%
    Je zwei Reispapierblätter kurz mit Wasser anfeuchten und direkt aufeinanderlegen.

    \step%
    Die doppelten Reispapierblätter in Quadrate schneiden und mit etwas Abstand auf das Blech oder Gitter legen.
    
    \step%
    Im Ofen 5--8 Minuten backen, bis die Stücke aufpuffen und kleine Kissen entstehen. Dabei gut beobachten, weil Reispapier ziemlich schnell von perfekt zu verbrannt wechseln kann.
    
    \step%
    Die Pillows vollständig auskühlen lassen. Danach mit einem heißen Metallspieß, einer Messerspitze oder vorsichtig mit einer Schere ein kleines Loch hineinmachen.
    
    \step%
    Für die Guacamole Avocado zerdrücken.
    Avocado mit Tomate, Zwiebel, Limettensaft, Salz, Pfeffer und Chiliflocken vermengen. Optional etwas Koriander oder Petersilie unterheben.
    
    \step%
    Guacamole in einen Spritzbeutel oder Gefrierbeutel mit abgeschnittener Ecke geben und vorsichtig in die Reispapier-Pillows füllen.
    
    \step%
    Direkt servieren, solange die Pillows noch knusprig sind.
}

\end{recipe}